% !TeX root = ../main.tex
% !TeX spellcheck = de_DE
\begin{otherlanguage}{ngerman}
\addchap{Kurzzusammenfassung}

% TODO: German version of the abstract, with the same structure and content.

\todo[inline]{\textbf{Struktur, nicht Format.}
  Die Bezeichnungen \emph{Kontext}, \emph{Problem}, \emph{Ziel}, \emph{Methode}, \emph{Ergebnisse} und \emph{Fazit} zeigen die Struktur, der die Kurzzusammenfassung folgen sollte.
  Sie müssen sie nicht als explizite Bezeichnungen verwenden: Sie können die Kurzzusammenfassung auch als Fließtext schreiben, solange sie diese Teile in dieser Reihenfolge abdeckt.}

\textit{Kontext.}
In welchem Bereich der Informatik ist Ihre Arbeit angesiedelt, und warum ist er relevant?

\textit{Problem.}
Welches Problem gehen Sie an, und warum wurde es bisher nicht zufriedenstellend gelöst?

\textit{Ziel.}
Was ist das Ziel Ihrer Arbeit, d.\,h.\ Ihre Forschungsfrage?

\textit{Methode.}
Wie sind Sie vorgegangen?
Was haben Sie entwickelt, gemessen oder untersucht?

\textit{Ergebnisse.}
Was sind Ihre wichtigsten Ergebnisse?

\textit{Fazit.}
Was bedeuten die Ergebnisse, und wer profitiert davon?
\end{otherlanguage}
